\section{Dinámica de Langevin: muestreo basado en Score}

\subsection{Idea central}

Hemos establecido anteriormente el concepto de \texttt{score function} y visto cómo codifica información rica sobre la distribución de probabilidad. La pregunta clave ahora es: si conocemos el \texttt{score function} de una distribución, ¿podemos muestrear de ella?

La respuesta es afirmativa; la dinámica de Langevin (\textbf{Langevin dynamics}) es precisamente el método para lograr este objetivo. Proviene de la física estadística y su idea central es que, mediante un proceso iterativo guiado por el \texttt{score function}, se puede partir de una \textbf{distribución inicial arbitraria} y converger finalmente a la distribución objetivo.

\begin{importantbox}{Proceso de muestreo de Langevin Dynamics}
Dada la función score del \texttt{score function} $\nabla_x \log q(x)$ de la distribución objetivo $q(x)$, la fórmula de actualización iterativa de la dinámica de Langevin es:
$$x_{k+1} = x_k + \eta\, \nabla_x \log q(x_k) + \sqrt{2\eta}\, z_k, \quad z_k \sim \mathcal{N}(0, I)$$
donde $\eta > 0$ es el parámetro de paso. Cuando $\eta \to 0$ y el número de iteraciones $K \to \infty$, la distribución de $x_K$ converge a $q(x)$.
\end{importantbox}

\begin{knowledgebox}{Los tres componentes de Langevin Dynamics}
La fórmula de actualización de la dinámica de Langevin consta de tres partes:
\begin{enumerate}
    \item \textbf{Posición actual} $x_k$: el estado actual de la partícula.
    \item \textbf{Guía por Score} $\eta\, \nabla_x \log q(x_k)$: la fuerza determinista que empuja la partícula hacia regiones de alta densidad de probabilidad.
    \item \textbf{Ruido aleatorio} $\sqrt{2\eta}\, z_k$: evita que la partícula se quede atrapada en valores locales óptimos y garantiza una exploración adecuada de la distribución objetivo.
\end{enumerate}
Si solo hubiera la guía por score sin inyección de ruido, la partícula convergería únicamente al punto de máximo de la densidad de probabilidad (modo), sin formar la distribución correcta. La existencia del término de ruido asegura la diversidad de las muestras finales.
\end{knowledgebox}

\subsection{Propiedades clave de Langevin Dynamics}

La dinámica de Langevin tiene una propiedad extraordinaria: \textbf{no requiere conocer} la función de densidad de probabilidad $q(x)$ en sí misma, sino únicamente el \texttt{score function} $\nabla_x \log q(x)$.

\begin{importantbox}{El Score Function basta para describir completamente la distribución}
El \texttt{score function} $\nabla_x \log q(x)$ determina de manera única la distribución de probabilidad $q(x)$ (hasta una constante de normalización), al cumplir ciertas condiciones de regularidad. Esto se debe a que:
\begin{itemize}
    \item $\nabla_x \log q(x) = \frac{\nabla_x q(x)}{q(x)}$, lo cual contiene información sobre el gradiente local de $q(x)$.
    \item Mediante la dinámica de Langevin, podemos recuperar el proceso de muestreo a partir del \texttt{score function}.
    \item Esto también implica que: incluso sin conocer la constante de normalización de la función de densidad de probabilidad (que suele ser difícil de calcular en espacios de alta dimensión), basta con conocer el \texttt{score function} para realizar el muestreo.
\end{itemize}
\end{importantbox}

En los modelos probabilísticos tradicionales, muestrear de distribuciones complejas requiere habitualmente calcular la inversa de la función de distribución acumulada (CDF), lo cual es prácticamente inviable en espacios de alta dimensión. La dinámica de Langevin ofrece una alternativa: basta con calcular el gradiente de la densidad de probabilidad logarítmica respecto a los puntos de datos y, mediante un proceso iterativo, se puede muestrear.

\begin{warningbox}{Limitaciones prácticas de Langevin Dynamics}
Las garantías teóricas de convergencia requieren que el paso $\eta \to 0$ y que las iteraciones sean infinitas; en la práctica, ambas condiciones no pueden cumplirse. El paso finito y el número finito de iteraciones introducen errores de aproximación. Además, en distribuciones multimodales, la dinámica de Langevin puede tardar mucho tiempo en saltar entre diferentes modos (problema de mezclado).
\end{warningbox}

\subsection{Resumen de la sección}

La dinámica de Langevin es un método de muestreo basado en el \texttt{score function}. A través del proceso iterativo de ``guía por score + inyección de ruido aleatorio'', puede partir de cualquier distribución inicial y converger a la distribución objetivo. La ventaja central de este método radica en que permite realizar el muestreo conociendo únicamente el \texttt{score function} y no la función de densidad de probabilidad en sí misma, evitando así el problema de calcular constantes de normalización difíciles en espacios de alta dimensión.


